\documentclass[11pt,a4paper]{article}
\usepackage[T1]{fontenc}
\usepackage[utf8]{inputenc}
\usepackage[ngerman]{babel}
\usepackage[margin=2.5cm]{geometry}
\usepackage[hidelinks]{hyperref}
\usepackage{array}
\setlength{\parindent}{0pt}
\setlength{\emergencystretch}{2em}
\renewcommand{\arraystretch}{1.25}
\begin{document}
\shorthandoff{"}
\begin{center}
  {\Large\textbf{Webtechnologien -- Übungsblatt 2}}
\end{center}
\section*{Aufgabe 2: Number-Guessing Game}
\textbf{Aufgabe:} Ein Zahlenratespiel mit \texttt{prompt}, \texttt{alert}
und \texttt{confirm} in reinem JavaScript umsetzen. Die HTML-Datei hat
einen leeren \texttt{body} und bindet nur die JavaScript-Datei ein.

\medskip
\textbf{Ergebnis:}
\begin{itemize}
  \item \texttt{index.html}: leerer \texttt{body}, \texttt{game.js} wird im
    \texttt{head} mit \texttt{defer} eingebunden.
  \item \texttt{MAX\_ATTEMPTS = 3} ist die Konstante für die Versuche,
    \texttt{keepPlaying} speichert, ob weitergespielt wird.
  \item Spielschleife: \texttt{while (keepPlaying)} erzeugt eine Zahl mit
    \texttt{Math.floor(Math.random() * 10) + 1}, ruft
    \texttt{playNumberGame} auf und fragt danach mit \texttt{confirm} nach
    einer weiteren Runde. Bei \glqq Abbrechen\grqq{} erscheint
    \glqq Thank you for playing!\grqq.
  \item Zufallszahl: \texttt{Math.random()} liefert $[0,1)$, mal 10 ergibt
    $[0,10)$, \texttt{Math.floor} ergibt 0 bis 9, plus 1 ergibt 1 bis 10.
    Beispiel: $0{,}42 \cdot 10 = 4{,}2 \rightarrow 4 \rightarrow 5$.
  \item \texttt{playNumberGame(secretNumber)}: speichert alle Tipps im Array
    \texttt{guesses} und läuft über höchstens drei Versuche.
    Fehlertext $\rightarrow$ \texttt{alert} mit der Meldung;
    richtig $\rightarrow$ Meldung mit Anzahl der Versuche,
    \texttt{return true}; sonst \glqq Too high\grqq{} bzw.
    \glqq Too low\grqq{} mit den bisherigen Versuchen und danach die
    bisherigen Tipps. Ohne Versuche: richtige Zahl anzeigen,
    \texttt{return false}.
  \item \texttt{getUserInput()}: wandelt die Eingabe mit \texttt{Number()}
    um. Im \texttt{try}-Block wird bei \texttt{NaN}, Werten unter 1 oder
    über 10 ein \texttt{Error} geworfen; \texttt{catch} gibt
    \texttt{error.message} zurück. Das Spiel erkennt den Fehlertext an
    \texttt{typeof guess !== "number"}.
\end{itemize}

\begin{center}
\begin{tabular}{|>{\raggedright\arraybackslash}p{3.6cm}|>{\raggedright\arraybackslash}p{4.2cm}|>{\raggedright\arraybackslash}p{5.6cm}|}
\hline
\textbf{Eingabe} & \textbf{Wert nach \texttt{Number()}} & \textbf{Ergebnis} \\ \hline
\texttt{3} & \texttt{3} & gültig \\ \hline
\texttt{11} & \texttt{11} & ungültig, größer als 10 \\ \hline
\texttt{abc} & \texttt{NaN} & ungültig \\ \hline
leere Eingabe & \texttt{0} & ungültig, kleiner als 1 \\ \hline
Abbrechen (\texttt{null}) & \texttt{0} & ungültig, kleiner als 1 \\ \hline
\end{tabular}
\end{center}

\textbf{Hinweis:} Die Angabe spricht bei der Spielentscheidung von einer
Konstante. \texttt{keepPlaying} ist trotzdem mit \texttt{let} deklariert,
weil der Wert nach jeder Runde neu gesetzt wird; eine \texttt{const} kann
nicht neu zugewiesen werden. Eine ungültige Eingabe zählt als Versuch.

\medskip
\textbf{Test:} In Chrome, \texttt{index.html} direkt geöffnet.
Geheimzahl jeweils 5:
\begin{center}
\begin{tabular}{|>{\raggedright\arraybackslash}p{3.6cm}|>{\raggedright\arraybackslash}p{10.1cm}|}
\hline
\textbf{Eingaben} & \textbf{Meldungen} \\ \hline
\texttt{3}, \texttt{5} & \glqq Too low! Attempts made: 1\grqq,
  \glqq Your previous guesses: 3\grqq,
  \glqq Correct! You needed 2 attempt(s).\grqq \\ \hline
\texttt{9}, \texttt{abc}, Abbrechen & \glqq Too high! Attempts made: 1\grqq,
  \glqq Your previous guesses: 9\grqq,
  zweimal \glqq Invalid input. Please enter a number between 1 and 10.\grqq,
  \glqq No attempts left. The correct number was 5.\grqq \\ \hline
\end{tabular}
\end{center}
Nach Abbrechen bei der Frage \glqq Do you want to play again?\grqq{}
erscheint jeweils \glqq Thank you for playing!\grqq.

\medskip
{\small \textbf{Hilfsmittel:} Claude Code (Claude Opus 5.5) für Erstellung, Prüfung und Dokumentation.}
\end{document}
